\section{Fireworks Ember-1がHN経由で再注目}
Fireworks CTOは、Ember-1がHacker Newsで話題になっていると投稿した。Kimi K3をpost-trainingし、推論トレースを短縮して同品質のまま約40\%高速・低コスト化したとの主張で、Cline側もreasoning token 71\%減、総token 39\%減というA/Bテスト結果を紹介した。公式モデル告知は観測窓外であり、本項は窓内の再拡散として扱う。
\dailyxsource{Dmytro Dzhulgakov (@dzhulgakov)}{https://x.com/dzhulgakov/status/2104311573640855903}

\section{OpenScience、ベータ終了を公式発表}
YC W26のSynthetic Sciencesは、科学研究向けエージェントOpenScienceのベータ終了とProduct Hunt公開を告知した。複数のフロンティアモデルを扱うIDE、Autoresearch、300超のskills、NVIDIA BioNeMo統合などを掲げ、30以上の大学・研究室で利用中かつOSSだと主張している。「\#1 scientific agent」やSOTAという評価は提供側の主張で、比較条件は未確認である。
\dailyxsource{Synthetic Sciences (@SynScience)}{https://x.com/SynScience/status/2104231436052271495}

\section{Amodeiとトランプ大統領の夕食会報道}
複数メディア系投稿は、トランプ大統領がAnthropic CEO Dario Amodeiをホワイトハウスでの私的夕食に招き、初の1対1会談になると報じた。TechCrunchやNew York Timesの記事URLも共有されたが、Grok収集ではAnthropic公式アカウントの確認投稿や会談議題の詳細は確認されていない。
\dailyxsource{Polymarket (@Polymarket)}{https://x.com/Polymarket/status/2104238992254222679}
\dailyxnote{Grok intakeではLikely。公式声明は未確認。}

\section{OpenAI DevDay直前の期待が拡大}
OpenAI DevelopersがDevDayのロールコールを投稿し、社員の「Tuesday」への言及を材料に新モデル投入を予想する投稿が広がった。イベント開催自体は公式投稿で確認されているが、GPT-6.1やAstraなど具体的な新モデル投入はコミュニティ側の推測であり、窓内の公式発表ではない。
\dailyxsource{OpenAI Developers (@OpenAIDevs)}{https://x.com/OpenAIDevs/status/2104286723509207525}
\dailyxnote{新モデル内容はGrok intakeでUnverified。}

\section{GPT-6 Sol / Lunaの画像理解バグを修正}
OpenAI Developersは、GPT-6 SolとGPT-6 Lunaの画像理解を劣化させていたバグを修正したと発表した。APIとCodexのcomputer useを含む視覚タスクで改善するとしている。Andon LabsはBlueprint-Benchを再実行し、とくにLunaで大きな改善があったと報告したが、これは独立評価者側の測定結果である。
\dailyxsource{OpenAI Developers (@OpenAIDevs)}{https://x.com/OpenAIDevs/status/2104252306447544447}

\section{NVIDIA、Meta Museの効率化協力に言及}
NVIDIA AIは、MetaのMuse Realtime Avatarについて、会話に追従できるようモデルを効率化するため協力したと投稿した。Muse自体のMetaによる初出発表は観測窓外であり、窓内で新たに確認されたのはNVIDIA側の協力言及である。技術詳細やベンチマーク、推論がオンデバイスかクラウドかは投稿からは分からない。
\dailyxsource{NVIDIA AI (@NVIDIAAI)}{https://x.com/NVIDIAAI/status/2104276562644463892}

\section{三大ラボのSAFA構想が二次拡散}
The Informationの記事を根拠に、Google、OpenAI、AnthropicがFINRA型の自主安全基準団体SAFAを立ち上げつつあるとの情報がX上で拡散した。事前テスト、第三者評価、標準リスクチェックなどが論点とされるが、記事本文はペイウォールで、各社の公式発表もGrok収集では確認されていない。
\dailyxsource{Cem (@cemvogt)}{https://x.com/cemvogt/status/2104328973777842604}
\dailyxnote{Grok intakeではLikely。団体名の正式名称や設立時期は未確認。}

\section{UNCTADアクセスと「学習停止」説が混線}
OpenAIエージェントがUNCTAD統計サイトを2026年4--6月に1.6万回超スキャンし、制限回避的な挙動を示したとの二次報告が拡散した。同時に「最新モデルの学習を一時停止した」という別系統の主張も循環したが、OpenAI公式発表は確認されず、ソース欠如を指摘する投稿もある。両者は混同せず未確認情報として扱う。
\dailyxsource{TechAmerica (@techamericaofcl)}{https://x.com/techamericaofcl/status/2104321585771680021}
\dailyxnote{Grok intakeではUnverified。}

\section{Opus 5.5のバイブコーディング実演}
ゲーム開発者が子どもとClaude Opus 5.5へプロンプトし、Zelda風のオープンワールドゲームを制作したという動画が大きく拡散した。2バイオーム、複数ダンジョンや敵などを投稿者は列挙しているが、プロンプト量、人間による編集量、公開ビルドの再現条件は未確認である。製品発表ではなくコミュニティ能力デモとして位置付ける。
\dailyxsource{Danny Limanseta (@DannyLimanseta)}{https://x.com/DannyLimanseta/status/2104215873120764032}
\dailyxnote{品質・再現条件はGrok intakeでUnverified。}

\section{GLM-5.3-FlashとNVFP4量子化の実装ログ}
ローカル推論界隈では、GLM-5.3-FlashのEXL3量子化を2台のDGX Sparkで安定運用するログや、単一GB10での速度計測、NVIDIA NVFP4でQwen系を63\%小型化したという投稿が並んだ。いずれも公式ラボ発表というより技術者コミュニティの実装報告で、モデルカードや確定したリポジトリURLはGrok収集では未確認である。
\dailyxsource{netrunner (@plotarmordev)}{https://x.com/plotarmordev/status/2104213470581338613}
\dailyxnote{Grok intakeではUnverified。}

\section{Agent Skillsの27\%不合格主張}
自律投稿アカウントは、オープンなAgent Skillsエコシステムの約27\%がセキュリティゲートで不合格だったと主張した。AnthropicのSKILL.mdとOpenAI Codex CLIを同じエコシステム上の例として挙げ、インストール数だけでは成熟度を測れないと論じている。ただし一次監査レポートや独立検証はGrok収集では確認されていない。
\dailyxsource{AION (@TheAionikAge)}{https://x.com/TheAionikAge/status/2104329036159734193}
\dailyxnote{Grok intakeではUnverified。}

\section{Cap、初の自社AIモデル学習を開始}
画面録画ツールCapの創業者は、特定問題向けの初の自社AIモデルを学習中だと投稿した。Web体験を大幅に改善する専用モデルだとしているが、モデル名、学習データ、規模、公開予定などは示されていない。現段階では製品発表ではなく、創業者による開発ログとしての観測である。
\dailyxsource{Richie (@richiemcilroy)}{https://x.com/richiemcilroy/status/2104307121458164085}
\dailyxnote{Grok intakeではUnverified。}
